\documentclass[11pt]{article}

\usepackage{amsmath}
\usepackage{amssymb}
\usepackage{amsthm}

\emergencystretch=2em

\newtheorem{theorem}{Theorem}
\newtheorem{lemma}[theorem]{Lemma}
\newtheorem{proposition}[theorem]{Proposition}
\newtheorem{corollary}[theorem]{Corollary}
\theoremstyle{definition}
\newtheorem{conjecture}[theorem]{Conjecture}
\newtheorem{definition}[theorem]{Definition}
\newtheorem{remark}[theorem]{Remark}

\newcommand{\R}{\mathbb{R}}
\newcommand{\N}{\mathbb{N}}
\newcommand{\norm}[1]{\left\lVert #1 \right\rVert}

\title{Refutation of Conjecture 00000004007:\\ the claimed optimal subadditivity inequality for $p$-variation is false}
\author{earthking11}
\date{13 September 2026}

\begin{document}
\maketitle

\begin{abstract}
We disprove Conjecture 00000004007.  The $p$-variation of a path is the
$p$-th root of the supremum over partitions of the sum of the $p$-th powers
of the increments.  Being a $p$-th root of a supremum, it is homogeneous of
degree one, $\norm{aX} = |a|\,\norm{X}$.  Taking $X = Y \neq 0$ gives
$\norm{X+Y} = 2\norm{X}$ on the left and $2^{1/q}\norm{X}$ on the right,
where $q = p/(p-1)$ is the conjugate exponent.  Since $q > 1$ for every
finite $p > 1$, we have $2 > 2^{1/q}$ and the conjectured inequality fails.
The same conclusion holds for the raw-supremum normalisation (without the
root), a fortiori, and for the limit case $p = 1$ ($q = \infty$).  The
failure is elementary; the inequality cannot be ``enlarged'' because it is
not valid in the first place.  The arithmetic core of the obstruction is
formalised in core Lean~4.
\end{abstract}

\section{The conjecture}

Let $X : [0,1] \to V$ be a path with values in a normed vector space $V$, and
let a \emph{partition} $\pi$ of $[0,1]$ be a finite increasing sequence
\[
  \pi : \quad 0 = t_0 < t_1 < \dots < t_k = 1 .
\]
The statement under consideration is the following.

\begin{conjecture}[00000004007]\label{conj}
With the $p$-variation defined as the $p$-th root of the supremum over
partitions of the sum of the $p$-th powers of the increments, and with $q$
the conjugate exponent of $p$, the inequality
\begin{equation}\label{eq:conj}
  \norm{X+Y} \le \bigl(\norm{X}^{q} + \norm{Y}^{q}\bigr)^{1/q}
\end{equation}
holds, and it cannot be enlarged.
\end{conjecture}

We show that \eqref{eq:conj} is false, so the conjecture is false.

\section{Definition and homogeneity}

\begin{definition}\label{def:pvar}
For $1 \le p < \infty$ the \emph{$p$-variation} of a path $X$ is
\begin{equation}\label{eq:def}
  \norm{X}
  := \left( \sup_{\pi} \sum_{i=1}^{k}
       \bigl\lVert X(t_i) - X(t_{i-1}) \bigr\rVert^{p} \right)^{1/p},
\end{equation}
the supremum being over all partitions $\pi = \{0 = t_0 < \dots < t_k = 1\}$.
The conjugate exponent of $p$ is $q := \dfrac{p}{p-1}$, characterised by
$\dfrac1p + \dfrac1q = 1$; thus $q \in (1,\infty)$ whenever
$p \in (1,\infty)$, and $q = \infty$ when $p = 1$.
\end{definition}

\begin{lemma}[Homogeneity]\label{lem:hom}
For every scalar $a$ and every path $X$,
\begin{equation}\label{eq:hom}
  \norm{aX} = |a| \, \norm{X}.
\end{equation}
\end{lemma}

\begin{proof}
The increments of $aX$ are $a\,(X(t_i) - X(t_{i-1}))$, so for each fixed
partition $\pi$,
\[
  \sum_{i=1}^{k} \bigl\lVert a\,(X(t_i) - X(t_{i-1})) \bigr\rVert^{p}
  = |a|^{p} \sum_{i=1}^{k} \bigl\lVert X(t_i) - X(t_{i-1}) \bigr\rVert^{p}.
\]
The factor $|a|^{p} \ge 0$ is independent of $\pi$, and for $c \ge 0$ one has
$\sup_\pi (c\, S_\pi) = c \, \sup_\pi S_\pi$.  Hence the $p$-th powers of the
suprema are related by $|a|^{p}$, and taking the $p$-th root gives
\[
  \norm{aX}
  = \bigl( |a|^{p}\, \norm{X}^{p} \bigr)^{1/p}
  = |a| \, \norm{X}. \qedhere
\]
\end{proof}

\begin{remark}
The homogeneity \eqref{eq:hom} is exactly the property that makes the
conjecture fail: it forces both sides of \eqref{eq:conj} to scale linearly
in the path, and the two linear coefficients disagree for $X = Y$.
\end{remark}

\section{Refutation: the case $X = Y$}

Let $X$ be any non-zero path and take $Y := X$.  Using Lemma~\ref{lem:hom},
the left-hand side of \eqref{eq:conj} is
\[
  \norm{X + Y} = \norm{2X} = 2 \norm{X},
\]
while the right-hand side is
\[
  \bigl(\norm{X}^{q} + \norm{Y}^{q}\bigr)^{1/q}
  = \bigl(2 \norm{X}^{q}\bigr)^{1/q}
  = 2^{1/q} \norm{X}.
\]
Since $X \neq 0$ we have $\norm{X} > 0$, so \eqref{eq:conj} at $X = Y$ is
equivalent to
\begin{equation}\label{eq:core}
  2 \le 2^{1/q}
  \quad\Longleftrightarrow\quad
  2^{q} \le 2 .
\end{equation}
But $q = p/(p-1) > 1$ for every finite $p > 1$, and $2^{q} > 2$ for every
$q > 1$.  Therefore \eqref{eq:core} is false, and the conjectured inequality
\eqref{eq:conj} fails.  It ``cannot be enlarged'' for the trivial reason that
it does not hold at all.

\begin{remark}[Elementary nature of the failure]
No estimate, tightness analysis, or special path is required.  The
counterexample is the simplest imaginable one, $X = Y \neq 0$; the failure is
a direct consequence of the degree-one homogeneity \eqref{eq:hom} and of
$q > 1$.  This is precisely why the stated inequality cannot be correct.
\end{remark}

\section{Explicit values for a two-point unit-step path}

Let $X$ be the two-point path with a single unit step, $X(0) = 0$,
$X(1) = 1$; there is only the trivial partition $\{0,1\}$, so
\[
  \norm{X} = \bigl( |1|^{p} \bigr)^{1/p} = 1 .
\]
For $Y = X$ the path $X+Y$ has a single step of size $2$, hence
\[
  \norm{X+Y} = \bigl( 2^{p} \bigr)^{1/p} = 2 ,
\]
independently of $p$, while the claimed right-hand side is
$\bigl(1^{q} + 1^{q}\bigr)^{1/q} = 2^{1/q}$.  The following table lists the
values; in every case the left-hand side strictly exceeds the right-hand
side.

\begin{center}
\begin{tabular}{c c c c c c}
\hline
$p$ & $q = \dfrac{p}{p-1}$ & $\norm{X}$ & $\norm{X+Y}$ & RHS $= 2^{1/q}$ & LHS $>$ RHS \\
\hline
$3/2$ & $3$      & $1$ & $2$ & $1.2599\ldots$ & yes \\
$2$   & $2$      & $1$ & $2$ & $1.4142\ldots$ & yes \\
$3$   & $3/2$    & $1$ & $2$ & $1.5874\ldots$ & yes \\
$4$   & $4/3$    & $1$ & $2$ & $1.6818\ldots$ & yes \\
$10$  & $10/9$   & $1$ & $2$ & $1.8661\ldots$ & yes \\
\hline
\end{tabular}
\end{center}

For $p = 2$ the requirement \eqref{eq:core} reads $4 \le 2$, and for $p = 3$
it reads $2^{3/2} \le 2$, i.e. $8 \le 4$; both are false.  As $p$ grows,
$q = p/(p-1) \downarrow 1$ and $2^{1/q} \uparrow 2$ from below, so the
margin $2 - 2^{1/q}$ shrinks but never vanishes for finite $p$: the
inequality fails for \emph{every} $p > 1$.

\section{The raw-supremum reading}

The conjecture specifies the $p$-th root in the definition
\eqref{eq:def}.  If instead one reads $\norm{\cdot}$ as the \emph{raw}
supremum, without taking the $p$-th root,
\begin{equation}\label{eq:raw}
  N(X) := \sup_{\pi} \sum_{i=1}^{k}
          \bigl\lVert X(t_i) - X(t_{i-1}) \bigr\rVert^{p},
\end{equation}
then $N$ is homogeneous of degree $p$, $N(aX) = |a|^{p} N(X)$.  At $X = Y$
with $\norm{X}$ denoting the raw quantity \eqref{eq:raw},
\[
  N(X+Y) = N(2X) = 2^{p} N(X),
\]
and the claimed right-hand side is
\[
  \bigl(2 N(X)^{q}\bigr)^{1/q} = 2^{1/q} N(X),
\]
so the inequality would require $2^{p} \le 2^{1/q}$.  For every $p \ge 1$ and
$q > 1$ we have $2^{p} \ge 2 > 2^{1/q}$, so the raw reading fails a
fortiori; for instance $p = 2$ gives $4$ on the left against
$2^{1/2} = 1.4142\ldots$ on the right.  Hence the refutation does not depend
on which of the two normalisations is intended.

\section{The limit case $p = 1$}

For $p = 1$ the quantity \eqref{eq:def} is the total variation and the
conjugate exponent is $q = \infty$, with
\[
  \bigl(\norm{X}^{q} + \norm{Y}^{q}\bigr)^{1/q}
  \xrightarrow[q \to \infty]{} \max\{\norm{X}, \norm{Y}\}.
\]
Taking $X = Y \neq 0$ gives
\[
  \norm{X+Y} = 2\norm{X} > \norm{X} = \max\{\norm{X}, \norm{Y}\},
\]
so \eqref{eq:conj} fails in this case as well.

\section{Conclusion}

Conjecture 00000004007 is false.  For every $p \in (1,\infty)$ the
conjugate exponent satisfies $q > 1$, and the single choice $X = Y \neq 0$,
together with the degree-one homogeneity of the $p$-variation, reduces the
claimed inequality to the false numerical statement $2^{q} \le 2$.  The
failure is elementary and occurs for all admissible $p$; the raw-supremum
reading and the limit case $p = 1$ fail as well.  There is no sense in which
the inequality can be ``enlarged'': it is simply not a true inequality.

\section{Formal verification}

The arithmetic core of the obstruction is formalised in core Lean~4
(\texttt{lean4/Main.lean}; \texttt{import Std}, no Mathlib, no \texttt{sorry},
no \texttt{axiom}, no \texttt{native\_decide}, no $\R$).  The formalised
statements are:
\begin{itemize}
  \item \texttt{q\_eq\_two\_fails} : $\lnot\,(2^{2} \le 2)$, i.e. $\lnot\,(4 \le 2)$
        --- the case $p = 2$ ($q = 2$);
  \item \texttt{q\_eq\_two} : $2^{2} = 4$;
  \item \texttt{q\_eq\_three\_halves\_fails} : $\lnot\,(2^{3} \le 2^{2})$,
        i.e. $\lnot\,(8 \le 4)$ --- the case $p = 3$ ($q = 3/2$, since
        $2^{3/2} \le 2 \iff 2^{3} \le 2^{2}$);
  \item \texttt{pow\_two\_gt\_two} : $\forall q : \N,\ 1 < q \to 2 < 2^{q}$,
        the general arithmetic fact behind \eqref{eq:core};
  \item \texttt{conjecture\_00000004007\_false}, collecting the two instances
        and the general fact; the accompanying comment records that
        homogeneity reduces the conjecture to exactly these false
        inequalities.
\end{itemize}
\texttt{lean4/Check.lean} runs \texttt{\#print axioms} for each theorem: the
\texttt{decide}-based theorems depend on no axioms, and
\texttt{pow\_two\_gt\_two} and the final theorem depend only on the core
axioms \texttt{propext} and \texttt{Quot.sound}; there is no \texttt{sorryAx}
and no \texttt{ofReduceBool}.  The real-valued $p$-variation itself, the
homogeneity Lemma~\ref{lem:hom} and the reduction of the conjecture to
\eqref{eq:core} are proved here and checked with exact rational arithmetic in
\texttt{reproduce.py}; they are not formalised in Lean (core Lean has no
$\R$).  See \texttt{lean4/README.md} for the precise scope.

\end{document}
